\documentclass[a4paper,12pt]{article}
\usepackage[T2A]{fontenc}
\usepackage[utf8]{inputenc}
\usepackage[russian]{babel}
\usepackage{amsmath,amssymb,mathtools}
\usepackage{cmbright}
\renewcommand{\familydefault}{\sfdefault}
\usepackage{icomma}
\usepackage{geometry}
\geometry{left=24mm,right=24mm,top=23mm,bottom=24mm}
\usepackage{microtype}
\usepackage{xcolor}
\usepackage{tikz}
\usetikzlibrary{calc,arrows.meta,positioning,shapes.geometric}
\usepackage{booktabs,tabularx}
\usepackage{enumitem}
\usepackage{parskip}
\usepackage{fancyhdr}
\usepackage{setspace}
\usepackage{hyperref}
\hypersetup{hidelinks}

\definecolor{accent}{HTML}{176B70}
\definecolor{darkaccent}{HTML}{194F55}
\definecolor{textgray}{HTML}{303638}
\color{textgray}
\onehalfspacing
\setlength{\headheight}{16pt}
\pagestyle{fancy}
\fancyhf{}
\fancyhead[L]{\small Ксения Васильченко}
\fancyhead[R]{\small 13.10.26}
\fancyfoot[C]{\small\thepage}
\renewcommand{\headrulewidth}{0.35pt}
\setlist[enumerate]{leftmargin=*,itemsep=.65em,topsep=.3em}
\newcommand{\sectionline}[1]{\section*{\color{darkaccent}#1}}
\newcommand{\learn}[1]{\par\noindent$\square$\enspace #1\par}
\newcommand{\keyidea}[1]{\par\noindent\textcolor{darkaccent}{\textbf{Запомните.}}\enspace #1\par}
\newcommand{\mfrac}[3]{#1\,\frac{#2}{#3}}

\begin{document}
\begin{titlepage}
\thispagestyle{empty}
\centering
\vspace*{32mm}
{\color{accent}\large МАТЕМАТИКА \textbullet\ 6 КЛАСС\par}
\vspace{18mm}
{\color{darkaccent}\Huge\bfseries Чек-лист\par}
\vspace{8mm}
{\LARGE Уравнения с дробями\par}
\vspace{3mm}
{\Large Вычитание смешанных чисел\par}
\vspace{3mm}
{\Large и действия со скобками\par}
\vspace{20mm}
{\large 13 октября 2026 года\par}
\vfill
{\large Лёвин Артём Александрович\par}
\vspace{2mm}
{\normalsize эксклюзивно для Ксении Васильченко\par}
\vspace{33mm}
\begin{tikzpicture}[remember picture,overlay]
\draw[accent!58,line width=1.25pt]
  ([xshift=25mm,yshift=20mm]current page.south west)
  .. controls ([xshift=79mm,yshift=42mm]current page.south west)
           and ([xshift=-77mm,yshift=4mm]current page.south east)
  .. ([xshift=-25mm,yshift=20mm]current page.south east);
\end{tikzpicture}
\end{titlepage}

\sectionline{1. Уравнения с обыкновенными дробями}
\learn{Вы находите неизвестное слагаемое, уменьшаемое или вычитаемое.}
\learn{Вы приводите дроби к общему знаменателю, сокращаете ответ и при необходимости выделяете целую часть.}

\begin{center}
\begin{tabular}{@{}ll@{}}
\toprule
\textbf{Вид уравнения} & \textbf{Как найти $x$}\\
\midrule
$x-a=b$ & $x=b+a$\\
$x+a=b$ & $x=b-a$\\
$a-x=b$ & $x=a-b$\\
\bottomrule
\end{tabular}
\end{center}

\textbf{Пример 1.} Найдём неизвестное уменьшаемое:
\[
 x-\frac7{20}=\frac45,\qquad
 x=\frac45+\frac7{20}=\frac{16}{20}+\frac7{20}
  =\frac{23}{20}=1\,\frac3{20}=1{,}15.
\]
Дробь $\frac{23}{20}$ неправильная: $23=20\cdot1+3$.
Десятичная запись получается делением $23$ на $20$.

\textbf{Пример 2.} Найдём неизвестное слагаемое:
\[
 x+\frac43=\frac{21}{4},\qquad
 x=\frac{21}{4}-\frac43
  =\frac{63}{12}-\frac{16}{12}
  =\frac{47}{12}=3\,\frac{11}{12}.
\]

\textbf{Пример 3.} Найдём неизвестное вычитаемое:
\[
 \frac{11}{12}-x=\frac1{18},\qquad
 x=\frac{11}{12}-\frac1{18}
 =\frac{33}{36}-\frac2{36}=\frac{31}{36}.
\]

\keyidea{Для сложения и вычитания дробей сначала найдите общий знаменатель. Если числитель и знаменатель имеют общий делитель больше $1$, сократите дробь. В частности, $\frac{31}{36}$ уже несократима: $\text{НОД}(31,36)=1$.}

\clearpage
\sectionline{2. Вычитание смешанных чисел}
\learn{Вы представляете смешанное число как сумму целой и дробной частей.}
\learn{Вы отдельно вычитаете целые и дробные части, если дробная часть уменьшаемого не меньше дробной части вычитаемого.}

\textbf{Пример 4.} В дробной части занимать единицу не требуется:
\begin{align*}
 4\,\frac56-2\,\frac14
 &=(4-2)+\left(\frac56-\frac14\right)\\
 &=2+\left(\frac{10}{12}-\frac3{12}\right)
 =2\,\frac7{12}.
\end{align*}

\keyidea{Разность смешанных чисел удобно считать по целой и дробной частям. Знак минус перед вторым смешанным числом относится \emph{и к его целой, и к его дробной части}.}

\sectionline{3. Если дробной части не хватает}
\learn{Вы умеете «занять» одну единицу из целой части.}
\learn{Вы понимаете, почему после такого преобразования значение числа не меняется.}

\textbf{Пример 5.} Вычислим $5\,\frac7{12}-1\,\frac58$.
Общий знаменатель дробных частей равен $24$:
\[
 5\,\frac7{12}-1\,\frac58
 =5\,\frac{14}{24}-1\,\frac{15}{24}.
\]
Поскольку $14<15$, возьмём одну единицу из числа $5$ и представим её как $\frac{24}{24}$:
\[
 5\,\frac{14}{24}
 =4+\frac{24}{24}+\frac{14}{24}
 =4\,\frac{38}{24}.
\]
Теперь можно вычитать:
\[
 4\,\frac{38}{24}-1\,\frac{15}{24}
 =(4-1)+\frac{38-15}{24}
 =\boxed{3\,\frac{23}{24}}.
\]

\keyidea{Если числитель дробной части уменьшаемого меньше соответствующего числителя вычитаемого, уменьшите целую часть на $1$, а к дробной части прибавьте $\frac{d}{d}$, где $d$ --- общий знаменатель.}

\clearpage
\thispagestyle{plain}
\sectionline{Алгоритм: как вычитать смешанные числа}
\vspace{5mm}
\begin{center}
\begin{tikzpicture}[
  >=Stealth,
  every node/.style={align=center,font=\small},
  process/.style={draw=accent,rounded corners=3pt,text width=9.7cm,minimum height=1.26cm,inner sep=5pt},
  decision/.style={draw=darkaccent,diamond,aspect=3.0,text width=5.3cm,inner sep=3pt},
  terminal/.style={draw=darkaccent,ellipse,minimum width=4.4cm,minimum height=1.0cm},
  arrow/.style={->,thick,draw=darkaccent},
  node distance=10mm and 11mm
]
\node[terminal] (start) {Начало};
\node[process,below=of start] (common) {Приведите дробные части обоих чисел к общему знаменателю $d$};
\node[decision,below=14mm of common] (compare) {Числитель первой\\дробной части не меньше\\второго числителя?};
\node[process,below=22mm of compare] (subtract) {Вычтите целые части и дробные части по отдельности};
\node[process,right=6mm of compare,text width=3.5cm] (borrow) {«Займите» единицу:\\целую часть уменьшите на $1$, дробную увеличьте на $\frac{d}{d}$};
\node[process,below=of subtract] (reduce) {Сократите дробь, если это возможно, и запишите ответ};
\node[terminal,below=of reduce] (finish) {Конец};
\draw[arrow] (start)--(common);
\draw[arrow] (common)--(compare);
\draw[arrow] (compare)--node[left]{Да}(subtract);
\draw[arrow] (compare.east)--node[above]{Нет}(borrow.west);
\draw[arrow] (borrow.south) |- (subtract.east);
\draw[arrow] (subtract)--(reduce);
\draw[arrow] (reduce)--(finish);
\end{tikzpicture}
\end{center}

\clearpage
\sectionline{4. Действия со скобками}
\learn{Вы сначала вычисляете выражение внутри скобок.}
\learn{Вы можете представить единицу дробью с подходящим знаменателем.}
\learn{Вы разбиваете длинное выражение на понятные последовательные действия.}

\textbf{Пример 6.} Вычислим выражение с вложенным вычитанием:
\[
\frac16-\left(1-\frac{17}{18}\right).
\]
Сначала выполним действие в скобках:
\[
1-\frac{17}{18}=\frac{18}{18}-\frac{17}{18}=\frac1{18}.
\]
Затем выполним внешнее вычитание:
\[
\frac16-\frac1{18}=\frac3{18}-\frac1{18}=\frac2{18}=\boxed{\frac19}.
\]

\textbf{Пример 7.} Вычислим по действиям:
\[
7\,\frac{14}{25}-\left(3\,\frac8{15}+1\,\frac9{10}\right).
\]
\textbf{Действие 1.} Сложим числа в скобках:
\begin{align*}
3\,\frac8{15}+1\,\frac9{10}
 &=4+\frac{16}{30}+\frac{27}{30}\\
 &=4+\frac{43}{30}
 =\boxed{5\,\frac{13}{30}}.
\end{align*}
\textbf{Действие 2.} Вычтем полученный результат:
\begin{align*}
7\,\frac{14}{25}-5\,\frac{13}{30}
 &=7\,\frac{84}{150}-5\,\frac{65}{150}\\
 &=\boxed{2\,\frac{19}{150}}.
\end{align*}

\keyidea{При нескольких действиях соблюдайте их порядок: сначала скобки, затем оставшиеся действия. Запись «Действие 1», «Действие 2» помогает сохранить логику вычислений.}

\clearpage
\sectionline{5. Самостоятельная тренировка \textnormal{(Путь А)}}
\textbf{10 заданий.} Решайте последовательно, записывая основные преобразования. В уравнениях найдите $x$; в остальных заданиях вычислите значение выражения. Ответы дайте в виде несократимых дробей или смешанных чисел.

\begin{enumerate}[itemsep=1.20em]
\item $x-\dfrac5{12}=\dfrac7{18}$.
\item $x-\dfrac35=\dfrac14$.
\item $x+\dfrac56=\dfrac94$.
\item $\dfrac{13}{15}-x=\dfrac16$.
\item $6\,\dfrac34-2\,\dfrac16$.
\item $8\,\dfrac3{10}-2\,\dfrac7{15}$.
\item $5\,\dfrac14-2\,\dfrac56$.
\item $\dfrac14-\left(1-\dfrac78\right)$.
\item $6\,\dfrac25-\left(1\,\dfrac34+2\,\dfrac1{10}\right)$.
\item $9\,\dfrac7{12}-\left(4\,\dfrac5{18}+1\,\dfrac58\right)$.
\end{enumerate}

\vspace{9mm}
\textbf{Самопроверка способа:} в заданиях 6, 7 и 10 потребуется проверить, хватает ли дробной части уменьшаемого. В заданиях 8--10 сначала выполните действия в скобках.

\clearpage
\sectionline{6. Ответы для самопроверки}
\renewcommand{\arraystretch}{1.38}
\begin{tabularx}{\linewidth}{@{}p{12mm}X@{}}
\toprule
\textbf{№}&\textbf{Ответ и ключевой шаг}\\
\midrule
1&$x=\frac{29}{36}$; $\frac{15}{36}+\frac{14}{36}=\frac{29}{36}$.\\
2&$x=\frac{17}{20}$; $\frac{12}{20}+\frac5{20}=\frac{17}{20}$.\\
3&$x=1\,\frac5{12}$; $\frac{27}{12}-\frac{10}{12}=\frac{17}{12}$.\\
4&$x=\frac7{10}$; $\frac{26}{30}-\frac5{30}=\frac{21}{30}$.\\
5&$4\,\frac7{12}$; $\frac34-\frac16=\frac7{12}$.\\
6&$5\,\frac56$; $8\,\frac9{30}=7\,\frac{39}{30}$.\\
7&$2\,\frac5{12}$; $5\,\frac3{12}=4\,\frac{15}{12}$.\\
8&$\frac18$; в скобках $\frac18$, затем $\frac14-\frac18=\frac18$.\\
9&$2\,\frac{11}{20}$; в скобках $3\,\frac{17}{20}$.\\
10&$3\,\frac{49}{72}$; в скобках $5\,\frac{65}{72}$, перед вычитанием $9\,\frac{42}{72}=8\,\frac{114}{72}$.\\
\bottomrule
\end{tabularx}

\vspace{8mm}
\keyidea{После вычисления проверьте, разумен ли ответ. Например, результат вычитания положительных чисел должен быть меньше уменьшаемого.}
\end{document}
